\label{sec:sigmamethod}
An even-parity block cannot represent the odd-parity $\sigma$ state.
Ref.~\cite{Montes2017}, Sec.~X, reports that the naive endpoint-spin and
single-fermion candidates do not directly yield the desired homogeneous
odd-sector ground state. To check the production covariance state
independently at small sizes, we construct a full-vector representative
of $\sigma$ by a Hamiltonian projection initialized from the block vacuum.
This seed and projection are our combination of the published block
vacuum and a standard eigensolver, not a formula claimed in that paper.
Here ``block-seeded'' means seeded by a conformal block; the solver is
single-vector Lanczos, not a multi-vector block-Lanczos algorithm.
The steps below specify the numerical construction, including the part
performed by the Lanczos solver~\cite{ARPACK,ScipyEigsh}.

\begin{enumerate}
\item \textbf{Prepare a seed with the required symmetry.}
Starting from normalized $\ket{I_{\rm CB}}$, set
\begin{equation}
 \ket{v_\sigma}=\sum_{j=0}^{N-1}X_j\ket{I_{\rm CB}},\qquad
 \ket{q_1}=\frac{\ket{v_\sigma}}{\|v_\sigma\|_2}.
 \label{eq:sigmaseed}
\end{equation}
Since $PX_j=-X_jP$, the seed is odd. The sum is invariant under
translation, so the seed has zero momentum. It generally contains
higher-energy odd states as well as $\sigma$; applying the summed spin
operator is not itself the final projection.
\item \textbf{Build the odd-sector Hamiltonian action.}
Retain only bit strings $\bm m$ with $\sum_jm_j$ odd, leaving
$2^{N-1}$ basis vectors. Let $H_-$ denote $H$ restricted to this basis.
For a coefficient vector $f$, its action is
\begin{equation}
 (H_-f)_{\bm m}=
 -\left[\sum_j(-1)^{m_j}\right]f_{\bm m}
 -\sum_j f_{\bm m\oplus\bm e_j\oplus\bm e_{j+1}},
 \label{eq:oddaction}
\end{equation}
where $\bm e_j$ has one bit at site $j$, $\oplus$ is bitwise exclusive-or,
and site indices are periodic. The first term is diagonal; the second
flips neighboring bits. No full $2^{N-1}$ by $2^{N-1}$ matrix is stored.
Only parity is explicitly imposed in this basis; translation symmetry
is inherited from the seed and preserved by the Hamiltonian.
\item \textbf{Generate a Krylov space.}
After $r$ basis vectors, the search space is
\begin{equation}
 \mathcal K_r=\operatorname{span}
 \{q_1,H_-q_1,\ldots,H_-^{r-1}q_1\}.
\end{equation}
Lanczos constructs orthonormal vectors $q_j$. In exact arithmetic,
with $q_0=0$ and $b_0=0$, the recurrence is
\begin{align}
 a_j&=\langle q_j|H_-|q_j\rangle,\notag\\
 \ket{w_j}&=H_-\ket{q_j}-a_j\ket{q_j}-b_{j-1}\ket{q_{j-1}},\notag\\
 b_j&=\|w_j\|_2,\qquad \ket{q_{j+1}}=\ket{w_j}/b_j.
 \label{eq:lanczos}
\end{align}
If $b_j=0$, the generated subspace is already invariant. In floating-point
arithmetic the solver also controls loss of orthogonality. The symbols
$a_j,b_j$ here are Lanczos coefficients, unrelated to the boson endpoint
labels or the leading fitting exponent $\alpha$.
\item \textbf{Select the lowest-energy combination.}
Form $(T_r)_{ij}=\langle q_i|H_-|q_j\rangle$. This small matrix is
tridiagonal, with diagonal $a_j$ and neighboring entries $b_j$.
Let $\theta_{\min}$ and $\bm y$ be its smallest eigenvalue and normalized
eigenvector. The Rayleigh--Ritz approximation is
\begin{equation}
 \ket{\sigma_N^{(r)}}=\sum_{j=1}^r y_j\ket{q_j}.
\end{equation}
It minimizes the energy within the current search space. Increasing the
space, or restarting while retaining the useful spectral information,
refines the approximation.
\item \textbf{Converge and measure the actual residual.}
The code calls \texttt{eigsh} with \texttt{k=1}, \texttt{which="SA"},
the seed as \texttt{v0}, and \texttt{tol=2e-12}. It uses implicitly
restarted Lanczos~\cite{ScipyEigsh}. The requested solver tolerance is
not an absolute bound on $\eta$: after convergence we independently
evaluate Eq.~\eqref{eq:eigenresidual} using the original Hamiltonian.
\item \textbf{Identify and validate the primary.}
Normalize the vector, restore zero entries on even bit strings, and
check $P=-1$, zero momentum, orthogonality to $I,\eps$, and the energy
$E_\sigma=-2\cot[\pi/(2N)]$. The vanishing large-$N$ gap has scaling
dimension $1/8$, identifying the $(1/16,1/16)$ primary.
Our explicit construction checks use
$e_P=\|(P+1)\sigma_N\|_2$ and
$e_{\mathsf T}=\|(\mathsf T-1)\sigma_N\|_2$ for the two symmetry errors.
\end{enumerate}

The spectral projection underlying these steps can also be written as
\begin{equation}
 \ket{\sigma_N}=\lim_{\tau\to\infty}
 \frac{e^{-\tau(H-E_*)}\ket{q_1}}
 {\|e^{-\tau(H-E_*)}\ket{q_1}\|_2}.
 \label{eq:projection}
\end{equation}
Here $\tau$ is imaginary time and $E_*$ is an arbitrary reference energy.
The normalized seed has an expansion into odd-sector eigenstates;
provided its $\sigma$ coefficient is nonzero, higher energies are
suppressed relative to the lowest one. Lanczos implements the search
without explicitly exponentiating $H$.
The measured squared seed overlap
$F_{\rm seed}=|\langle\sigma_N|q_1\rangle|^2$ is approximately
$0.986,0.983,0.981,0.979,0.977$ at $N=6,8,10,12,16$.
At $N=16$, the recorded projected-state residual is
$2.70\times10^{-12}$.
The complete size-by-size residual and independent-state checks appear
in Appendix~\ref{sec:ising-spin-validation}; the largest recorded residual is
$1.99\times10^{-11}$ at $N=12$. Thus the solver tolerance must not be
quoted as the achieved absolute residual.

This method obtains the required lattice primary using both a conformal
block and the Hamiltonian. A pure endpoint conformal-block formula for
the periodic $\sigma$ state remains unestablished here.
